\section{Summary}
\label{sec:conclusion:summary}

This report develops the design of an analytical and visualisation
platform for the Ho Chi Minh City public bus network.
Chapter~\ref{chapter:intro} motivates the work as the engineering
response to the absence of any analytical surface over the network's
two operational data streams, namely Automated Fare Collection (AFC)
and Global Positioning System (GPS) telemetry.
Chapter~\ref{chapter:related} surveys the four existing tools for the
network (a public route directory and three passenger routing
applications) and the peer-reviewed research that informs the design
(the Bangkok Quality-of-Service framework~\cite{srisakaokul2023}, the
bus-as-traffic-sensor construction~\cite{derevitskiy2016}, the
time-distance trajectory predictor~\cite{linzeng2014}, the
visualisation primitives of Liu~\cite{liu2022}, and the storytelling
principles of Knaflic~\cite{knaflic2015}); the chapter identifies the
gap that the present project addresses.

Chapter~\ref{chapter:foundations} introduces the theoretical and
technical foundations on which the proposed solution depends.
Chapter~\ref{chapter:solution} develops the proposed solution along
eight axes: problem analysis, dataset and exploratory analysis,
preprocessing, solution overview, data architecture, analytical
themes, visualisation technique selection, and interaction model.
Chapter~\ref{chapter:implementation} documents the Phase~One artefacts,
namely Layer~1 of the warehouse, the user-interface design mockup, and
the initial requirements evaluation.
Chapter~\ref{chapter:future} sets out the Phase~Two roadmap of nine
phases and approximately twenty-one working days, the deployment plan,
the user-evaluation methodology, and three extensions that follow from
the historical foundation.

The proposed solution composes a four-layer warehouse, six analytical
themes, a visualisation catalogue grounded in established primitives,
and an interaction model that combines Multiple Coordinated Views,
visual storytelling, and an entity-detail interaction pattern. Each
design choice is defended against at least one rejected alternative.
The composition responds to the analytical gap identified in
Chapter~\ref{chapter:related} and forms the basis for the Phase~Two
implementation.
